% se puede agregar la opción [english] para 
%  memorias o tesis en inglés (borrando el archivo .aux)
\documentclass{umemoria} 

\facultad{Facultad de Ciencias de la Salud}
\carrera{Carrera de Química Y Farmacia}
\tesis{Tesis para Optar al Grado de\\Químico Farmacéutico}
\author{Nombre Completo del Tesista} %Síntesis y Caracterización de Derivados de 1,3,5-triazina-Escuaramidas
\title{Título de la Tesis} %Síntesis y Caracterización de Derivados de 1,3,5-triazina-Escuaramidas
\lugar{Santiago -- Chile\\ SEDE SANTIAGO, CAMPUS PROVIDENCIA}
\anho{Marzo, 2023}

% puede haber varios profesores guía seperados por coma;
% pero si es una memoria, solo puede haber un profesor guía
\guia{Nombre Completo Guía} %Dr. Desmond A. MacLeod-Carey Castro

% puede haber varios profesores co-guía seperados por coma;
% pero si es una memoria, el profesor co-guía será el primer
% integrante de la comisión
%\coguia{Nombre Completo Co-Guía} % incluir en caso de co-guía de *tesis*
%\cotutela{Nombre Institución} % incluir en caso de cotutela

\comision{Nombre Completo Uno,Nombre Completo Dos,Nombre Completo Tres}

%\auspicio{Nombre Institución} % incluir en caso de recibir financiamiento

% tiene que ser el año en que se da el examen de título/grado (defensa)
%\anho{2021} % incluir solo para reemplazar el año actual

\usepackage{lipsum}
\usepackage[version=4]{mhchem}
\usepackage{xcolor}

\begin{document}

\frontmatter
\maketitle

\begin{resumen}
\lipsum[1-3]
\end{resumen}

% opcional: incluir para tesis en inglés;
%  en este caso hay que tener el resumen y abstract
%   en ambos idiomas
\begin{abstract}
\lipsum[1-3]
\end{abstract}

\begin{dedicatoria}
Una dedicatoria corta.
\end{dedicatoria}

\begin{thanks}
\lipsum[1-2]
\end{thanks}

\tableofcontents
\listoftables % opcional
\listoffigures % opcional

\mainmatter

\input{cap1-Intro.tex}
\input{cap2-Metod.tex}
\input{cap3-Result.tex}
\input{cap4-Conclu.tex}

% ver https://www.overleaf.com/learn/latex/Glossaries
% \input{glosario.tex} % opcional

\nocite{*}
\bibliographystyle{plain}
\bibliography{bibliografia}

% opcional ...
\begin{appendices}
\input{anexoA.tex}
\end{appendices}
\end{document}